% =====================================================================
% preambulo.tex
% Responsabilidad única: cargar paquetes y definir configuración/estilos
% reutilizables (colores de enlaces, entornos, listas). No contiene
% contenido propio del taller — si algo del "look" del documento debe
% cambiar, este es el único archivo que hay que tocar.
% =====================================================================

\usepackage[utf8]{inputenc}
\usepackage[T1]{fontenc}
\usepackage[spanish,es-noshorthands]{babel}
\usepackage[a4paper,margin=2.5cm]{geometry}

% --- Matemáticas ------------------------------------------------------
\usepackage{amsmath}
\usepackage{amssymb}
\usepackage{amsthm}
\usepackage{mathtools}

% --- Tablas y figuras ---------------------------------------------------
\usepackage{graphicx}
\usepackage{booktabs}
\usepackage{array}
\usepackage{multirow}
\usepackage{caption}
\usepackage{float}

% --- Listas con formato a., b., c. ------------------------------------
\usepackage{enumitem}

% --- Diagramas vectoriales -----------------------------------------------
\usepackage{tikz}
\usetikzlibrary{arrows.meta, positioning, calc}
% Misma paleta usada en las figuras generadas con matplotlib (Ejercicios
% 6-8), para que los diagramas TikZ se vean como parte del mismo documento.
\definecolor{classblue}{HTML}{4C72B0}
\definecolor{classred}{HTML}{C44E52}
\definecolor{classgreen}{HTML}{55A868}

% --- Color y enlaces ----------------------------------------------------
\usepackage{xcolor}
\definecolor{urosariored}{RGB}{255,11,38} % rojo exacto muestreado del escudo
\usepackage{hyperref}
\hypersetup{
    colorlinks = true,
    linkcolor  = urosariored!75!black,
    citecolor  = urosariored!75!black,
    urlcolor   = urosariored!75!black,
    pdftitle   = {Taller 1 -- KNN y Metricas de Clasificacion},
    pdfauthor  = {Over Haider Castrillon Valencia},
}

% --- Espaciado de tablas ------------------------------------------------
\renewcommand{\arraystretch}{1.25}

% --- Pie de página institucional -----------------------------------------
% Responsabilidad de este bloque: mostrar el escudo de la Universidad del
% Rosario, pequeño, en la esquina inferior derecha de cada página del
% cuerpo del documento (la portada define su propio pie por separado en
% portada.tex, ya que \begin{titlepage} usa página sin estilo).
\usepackage{fancyhdr}
\pagestyle{fancy}
\fancyhf{}
\renewcommand{\headrulewidth}{0pt}
\renewcommand{\footrulewidth}{0pt}
\fancyfoot[C]{\thepage}
\fancyfoot[R]{\includegraphics[height=0.42cm]{figuras/logo_urosario_icono.png}}

% --- Encabezados de subsección sin numerar, para incisos a-e -----------
\setlist[enumerate,1]{label=\alph*., leftmargin=1.6em}
